\chapter{Reduced-Order Gait Optimization}
\label{ch:mujoco}

This chapter addresses RQ1. A whole-robot model with quasi-steady link forces is built in MuJoCo
\cite{Todorov2012}, and its gait is optimized with CMA-ES \cite{Hansen2001,Hansen2016} under actuator and
attitude constraints. The model is cheap enough for thousands of evaluations but its coefficients are not
calibrated. Its purpose is to show which propulsion principle an optimizer selects, and how the answer depends on the physics included, not to predict speeds of BODY2.

\section{Model}
\label{sec:mj-model}

\subsection{Robot}

BODY2 could not be used directly: its CAD-derived model is an open kinematic tree in which the parallelogram
linkages are not closed, and its buoyancy distribution had not been measured. The optimization therefore uses
a simple quadruped model, \texttt{toy\_quad}, of similar size (\cref{fig:mj-model}, \cref{tab:mj-robot}). Each
leg has three hinge joints about the transverse axis: hip, knee and flipper wrist, twelve actuators in total.
The flipper, a thin plate of \qtyproduct{52 x 36}{\milli\metre} whose area of \SI{18.7}{\square\centi\metre} is close to that of the open fan paddle of BODY2, is called \emph{flipper} throughout to distinguish it from the fan paddle and the fluke of
the real robot: it can paddle or flap depending on the physics included, whereas the two propulsors of BODY2 are each built for one
principle. The dry mass is \SI{0.78}{\kilo\gram}; fully submerged, the robot would
displace \SI{1126}{\cubic\centi\metre}, so it floats at the surface like BODY2. The body length of
\SI{0.239}{\metre} is used to express speeds in body lengths per second.

\begin{figure}[t]
  \centering
  \includegraphics[width=0.42\textwidth]{fig_mj_model.png}
  \caption[MuJoCo model]{The MuJoCo model \texttt{toy\_quad} floating at the surface (rendering from the
    simulation). Each leg has a hip, a knee and a flipper joint.}
  \label{fig:mj-model}
\end{figure}

\begin{table}[t]
  \centering
  \caption[MuJoCo robot model]{The reduced-order robot model \texttt{toy\_quad}.}
  \label{tab:mj-robot}
  \small
  \begin{tabular}{@{}llll@{}}
    \toprule
    Part & Geometry & Size & Mass \\
    \midrule
    Trunk & box & \qtyproduct{0.20 x 0.10 x 0.05}{\metre} & \SI{0.600}{\kilo\gram} \\
    Upper leg ($\times4$) & capsule & $r=\SI{7}{\milli\metre}$, $l=\SI{55}{\milli\metre}$ & \SI{0.020}{\kilo\gram} \\
    Lower leg ($\times4$) & capsule & $r=\SI{6}{\milli\metre}$, $l=\SI{50}{\milli\metre}$ & \SI{0.015}{\kilo\gram} \\
    Flipper ($\times4$) & plate & \qtyproduct{8 x 52 x 36}{\milli\metre} & \SI{0.010}{\kilo\gram} \\
    \midrule
    \multicolumn{4}{@{}l}{Joint ranges: hip and knee $\pm\SI{1.3}{\radian}$, flipper $\pm\SI{1.6}{\radian}$;
      legs at $\pm\SI{0.085}{\metre}$ fore--aft, $\pm\SI{0.045}{\metre}$ lateral} \\
    \multicolumn{4}{@{}l}{Dry mass \SI{0.78}{\kilo\gram}; added mass \SI{0.28}{\kilo\gram}; body length
      \SI{0.239}{\metre}; time step \SI{2}{\milli\second}, implicit integrator} \\
    \bottomrule
  \end{tabular}
\end{table}

\subsection{Hydrodynamic Forces}

Each geometric body $k$ receives a quasi-steady force, scaled by its immersed fraction $f_k$ and applied at its
centre:
\begin{equation}
  \bm F_k = \rho g V_k f_k\,\bm e_z
  \;-\; \tfrac12\rho f_k \sum_{i=1}^{3} C_{D,i} A_i\, v_i|v_i|\,\bm e_i
  \;-\; c_v f_k\,\bm v \;+\; \bm L_k .
  \label{eq:mj-force}
\end{equation}
Here $\bm v$ is the velocity of the body centre relative to still water, $v_i$ its components along the
body axes $\bm e_i$, and $A_i$ the projected areas; $c_v$ is a small linear damping coefficient (0.05, 0.02 and \SI{0.10}{\newton\second\per\metre} for flipper, links and trunk) that matters only at very low speed. The drag is anisotropic: the flipper has
$C_D=(2.2,\,0.10,\,0.10)$ normal to and along its faces, the links $(0.5,0.5,0.25)$ and the trunk
$(0.9,1.6,1.8)$. This anisotropy is the only source of thrust in a drag-only model. The immersed fraction is
computed from the actual vertical extent of each body, so a flipper that breaks the surface loses its force.
Added mass $m_a=C_a\rho V$ ($C_a=1.0$ for the flipper, 0.3 for the links, 0.2 for the trunk) is added to the
body inertia together with a constant upward force $m_a g$ that cancels its weight. This isotropic
approximation is unconditionally stable, unlike an explicit $\dot{\bm v}$ term. The coefficients are
literature values of the same level as in beaver-robot models \cite{Chen2021,Chen2022}; they are not
calibrated.

\paragraph{Lift.} All terms of \eqref{eq:mj-force} except $\bm L_k$ oppose some component of the relative flow.
An optional post-stall flat-plate lift acts on the flippers,
\begin{equation}
  C_L(\alpha) = c_l \sin 2\alpha, \qquad |\bm L| = \tfrac12\rho\,C_L\,A_\text{face}|\bm v|^2 f,
  \label{eq:mj-lift}
\end{equation}
normal to $\bm v$ in the plane spanned by $\bm v$ and the plate normal, where $\alpha$ is the angle between
$\bm v$ and the plate. The value $c_l=1.1$ is a textbook flat-plate value without an aspect-ratio correction.
At small $\alpha$ the lift grows like $\alpha$ and the drag like $\alpha^2$, so a fast, shallow, feathered
stroke produces almost no thrust in the drag-only model but a large thrust with lift. With the lift switched
off, a lift-based gait cannot win, because the mechanism that would make it win is absent.

\paragraph{Thrust decomposition.} To identify the mechanism of a gait, the signed impulses of the lift terms and
of the drag terms (including the linear term) along the swimming direction are accumulated over the scoring
window. A gait with positive lift impulse and negative drag impulse is \emph{lift-propelled}: lift supplies the
forward impulse, and drag only opposes it. The decomposition was checked with a coasting test, in which the
initial momentum equals the sum of the impulses. To quantify how much of the propulsion lift provides, the
\emph{lift thrust share}
\begin{equation}
  s_L = \frac{I_L}{|I_\text{trunk}|}
  \label{eq:liftshare}
\end{equation}
relates the forward impulse $I_L$ of the lift terms to the drag impulse of the trunk; the drag on the limbs is
accumulated separately. A paddling stroke has $s_L\approx0$; a lift-propelled stroke has $s_L>1$, because lift must also
overcome the drag of the limbs themselves. The share can be imposed as a further limit, $s_L\le s_\text{max}$, which forces
drag-based strokes in a model that contains lift (\cref{sec:mj-liftshare}).

\subsection{Servo Model}
\label{sec:mj-servo}

The actuators are position servos with a stiffness of \SI{60}{\newton\metre\per\radian} and critical damping. Real hobby servos cannot move
arbitrarily fast. Each actuator is therefore given the torque--speed characteristic of a DC motor,
\begin{equation}
  \tau = \tau_s\,u - \frac{\tau_s}{\omega_\text{max}}\,\omega, \qquad |u|\le1,
  \label{eq:servo}
\end{equation}
with stall torque $\tau_s=\SI{3}{\newton\metre}$ and no-load speed
$\omega_\text{max}=\SI{400}{\degree\per\second}$. These values were chosen before the servos of BODY2 were identified; the PTK\,7465W has a lower stall torque and a higher no-load speed (\cref{sec:plat-robot}), and the consequences are discussed in \cref{sec:disc-limits}. The back-EMF term is a viscous damping
$b=\tau_s/\omega_\text{max}\approx\SI{0.43}{\newton\metre\second\per\radian}$, which is added to the joint
damping and integrated implicitly by MuJoCo; the result is independent of the time step. Two simpler
variants failed: limiting only the rate of the command let the legs be dragged back and then snap forward at
\SI{1500}{\degree\per\second}, and rewriting the torque limit at every step acted as an explicit damping that
oscillated at up to \SI{3000}{\degree\per\second} and changed the result six-fold with the time step.

Two power measures are reported. The optimization uses the mean mechanical power at the motor shafts,
$\sum_j\langle|\tau_{m,j}\,\omega_j|\rangle$ with the motor torque $\tau_m=\tau-b\,\omega$ and no recovery of
braking energy. The electrical power adds the copper loss of the same motor model,
$\sum_j\langle\max(\tau_{m,j}\omega_j,0)+\tau_{m,j}^2\,\omega_\text{max}/\tau_s\rangle$, which follows from
\eqref{eq:servo} without new parameters; it is computed after the optimization and used as a check. The cost of
transport in this chapter is power over speed in \si{\joule\per\metre}.

\section{Gait Parametrization}
\label{sec:mj-gait}

Every actuator $i$ follows the target angle
\begin{equation}
  q_i(t) = q_{0,i} + r(t)\sum_{h=1}^{H} A_{i,h}\,\sin\!\big(h\,\vartheta(t)+\varphi_{i,h}\big), \qquad H=2,
  \label{eq:gait}
\end{equation}
with offset $q_{0,i}$, amplitudes $A_{i,h}$, phases $\varphi_{i,h}$ and a ramp $r(t)$ from 0 to 1 over the
first \SI{1.5}{\second}. The duty warp $\vartheta(t)$ lets the power stroke be faster or slower than the recovery
while keeping the period: with the cycle progress $c=(ft)\bmod1$,
\begin{equation}
  \vartheta = \begin{cases}
    \pi\,c/\mathrm{DUTY}, & c<\mathrm{DUTY},\\
    \pi + \pi\,(c-\mathrm{DUTY})/(1-\mathrm{DUTY}), & c\ge\mathrm{DUTY}.
  \end{cases}
  \label{eq:duty-warp}
\end{equation}
The second harmonic, which is invariant under a phase shift of $\pi$, can express the feathering of the
flipper. It is not necessary for net thrust: a single harmonic already propels the model. Joints of the same
type share amplitude and offset. The search ranges are $f\in[0.2,2.5]$\,Hz, $\mathrm{DUTY}\in[0.25,0.75]$,
$A\in[0,0.9]$\,rad and $q_0\in[-0.5,0.5]$\,rad, normalized to $[0,1]$.

The coordination of the legs is either free, with an independent phase for every actuator (35 parameters),
or fixed to a \emph{gait family} with one phase per joint type (17 parameters). The families are defined by
the delay of each leg relative to the left fore leg FL, applied in time before the duty warp:
\begin{center}
  \small
  \begin{tabular}{@{}lll@{}}
    \toprule
    Family & Common name & Delays of FR, BL, BR (fraction of period) \\
    \midrule
    \texttt{diag} & trot & 0.5, 0.5, 0 \\
    \texttt{lr} & pace & 0.5, 0, 0.5 \\
    \texttt{fb} & bound & 0, 0.5, 0.5 \\
    \texttt{wave} & walk & 0.5, 0.75, 0.25 \\
    \texttt{inphase} & pronk & 0, 0, 0 \\
    \bottomrule
  \end{tabular}
\end{center}
To classify a free gait, the delays of the joint type with the largest first-harmonic amplitude are compared
with each family, and the family with the smallest mean circular distance is reported together with that
distance (0 = identical, 0.5 = opposite).

\section{Objective, Constraints and Optimizer}
\label{sec:mj-objective}

Each evaluation settles the robot for \SI{1.0}{\second}, ramps the gait in over \SI{1.5}{\second} and scores the
following \SI{8}{\second}, rounded up to whole stroke periods. The speed is the displacement along the initial
heading divided by the window length, so turning is penalized automatically. Two objectives are used: the
highest speed, or the lowest mean mechanical power subject to a minimum speed $v_\text{min}$. All settings are listed in \cref{tab:app-mujoco}.

Stability is imposed by limits, not by weights. With RMS values of heading error, roll and pitch and their
limits $\ell_k$, the penalty is
\begin{equation}
  p = \sum_k \frac{\max(0,\mathrm{RMS}_k-\ell_k)}{\ell_k}
    \;\big[+\,\text{surfacing and } v_\text{min} \text{ terms}\big], \qquad
  \text{fitness} = \begin{cases} \text{objective}, & p=0,\\ -1000-p, & p>0. \end{cases}
  \label{eq:feasibility}
\end{equation}
Any infeasible gait thus scores below every feasible one; speed cannot buy back a violation. An earlier
additive weighting of speed and attitude failed once the legs could move freely: speed dominated, and the best
gait rolled by \SI{42}{\degree} and yawed by \SI{44}{\degree}. Two limit sets are used: a moderate one,
heading/roll/pitch $10/10/\SI{15}{\degree}$, used in this chapter, and a strict one, $10/5/\SI{5}{\degree}$, used for the
coordination study in \cref{app:coord}. Where stated, the
fraction of time in which any limb is less than half immersed (surfacing) is limited to 5\,\%, because the
model has no free-surface physics.

The optimizer is CMA-ES with population 16 and initial step size 0.25 in the normalized space, run for 2500
evaluations unless stated otherwise, with several seeds per condition. Each result is a bound of what the model
allows: a lower bound for the highest speed and an upper bound for the least power. Values are given as mean
$\pm$ standard deviation over seeds.

\section{Results}
\label{sec:mj-results}

\subsection{Baseline and Servo Limit}

A hand-designed trot at \SI{1.2}{\hertz} swims at \SI{0.070}{\metre\per\second} (0.29 body lengths per second)
with \SI{26.7}{\degree} RMS pitch and servos saturated 85\,\% of the time. Without the servo model, the
optimizer found gaits with joint speeds of \SIrange{2700}{4200}{\degree\per\second}, about seven to ten times the servo limit. Replayed with the servo model, the best of them did not move forward
(\SI{-0.001}{\metre\per\second}). Re-optimized with the servo model, three free-phase searches reached
\SIrange{0.11}{0.25}{\metre\per\second}, all close to a bound. The servo limit therefore has to be part of the
model; all following results include it.

\subsection{Drag versus Lift at Equal Speed}

Comparing the fastest gaits of the two physics would be unfair, because power rises steeply with speed and the
fastest gaits differ in speed. The comparison is therefore made at equal speed: the least mechanical power
needed to reach $v_\text{min}\in\{0.10,0.20,0.30\}$\,m/s under moderate limits and at most 5\,\% surfacing,
over five families and two seeds (\cref{tab:mj-pareto}, \cref{fig:mj-pareto}).

\begin{table}[t]
  \centering
  \caption[Least power at equal speed, drag versus lift]{Least mean mechanical power to reach a required speed,
    moderate limits, best over five families and two seeds (family of the best solution and number of feasible searches out of
    ten in brackets). Left: drag-based strokes in a model without lift; middle: drag-based strokes forced by $s_L\le0.1$ in the
    model with lift (\cref{sec:mj-liftshare}); right: unconstrained strokes in the model with lift.}
  \label{tab:mj-pareto}
  \small
  \begin{tabular}{@{}cccc@{}}
    \toprule
    Required speed & Drag-only model & Lift model, $s_L\le0.1$ & Lift model, free \\
    \midrule
    \SI{0.10}{\metre\per\second} & \SI{0.23}{\watt} (walk) & \SI{0.42}{\watt} (bound; 2/10) & \SI{0.10}{\watt} (bound; 10/10) \\
    \SI{0.20}{\metre\per\second} & \SI{1.04}{\watt} (bound) & \SI{5.50}{\watt} (pronk; 1/10) & \SI{0.55}{\watt} (pronk; 8/10) \\
    \SI{0.30}{\metre\per\second} & \SI{3.10}{\watt} (trot) & none (0/10) & \SI{0.97}{\watt} (bound; 4/10) \\
    \bottomrule
  \end{tabular}
\end{table}

Three results are robust. First, whenever lift is in the model and not suppressed, every feasible optimum is
lift-propelled: in all 22 feasible lift optima of this experiment the lift impulse is positive
(\SIrange{0.46}{2.2}{\newton\second} on the front) and the drag impulse negative. The optimizer never paddles when it can
flap. Second, at equal speed
the lift-propelled gaits need about one half to one third of the power of the best drag-only gaits, at all three
speeds. Third, each optimum depends on its physics: with lift removed, the lift optima keep only 35--58\,\% of
their speed, and 26 of the 27 drag optima no longer satisfy the attitude limits when lift is added, because lift also
creates moments on the trunk (\cref{sec:mj-liftshare}). The size of the ratio is not reliable: the searches have not converged (a second seed lowered the
drag-only power at \SI{0.20}{\metre\per\second} from 2.99 to \SI{1.04}{\watt}), so both fronts are upper bounds. The ranking is the same with the electrical power, which includes the copper losses of the servos: the drag-only front needs 0.22, 1.14 and \SI{2.80}{\watt}, the lift front 0.10, 0.47 and \SI{0.86}{\watt}.

\begin{figure}[t]
  \centering
  \includegraphics[width=0.62\textwidth]{fig_mj_pareto}
  \caption[Least power at equal speed]{Least mechanical power to reach a required speed, moderate limits: drag-only model, model
    with lift, and drag-based strokes forced by $s_L\le0.1$ in the model with lift. Pale markers are the individual optima of each
    family and seed.}
  \label{fig:mj-pareto}
\end{figure}

\subsection{Drag-Based Strokes within the Lift Model}
\label{sec:mj-liftshare}

The comparison of \cref{tab:mj-pareto} between the drag-only model and the lift model is a comparison of two models as
much as of two propulsion principles. A direct test of RQ1 compares both principles within one model. For this, the
least-power searches of the lift model were repeated with the additional limit $s_L\le0.1$ from \eqref{eq:liftshare}, which
admits only strokes that obtain their thrust from drag, with the same speeds, limits, families and seeds as the unconstrained
searches (30 searches).

The results are shown in the middle column of \cref{tab:mj-pareto} and in \cref{fig:mj-pareto}. Only 3 of the 30 searches
found a feasible drag-based stroke; most of the others satisfy the share limit but not the required speed. On this small basis,
the drag-based strokes need 4.2 times the power of the lift-propelled ones at \SI{0.10}{\metre\per\second} (\SI{0.42}{} against
\SI{0.10}{\watt}) and 10 times at \SI{0.20}{\metre\per\second} (\SI{5.50}{} against \SI{0.55}{\watt}), and none reaches
\SI{0.30}{\metre\per\second}, the closest reaching \SI{0.26}{\metre\per\second}. Because the feasible region is narrow, part of the
factor may be due to search difficulty rather than physics; with two seeds and 2500 evaluations, the direction is reliable but the
value is not.

The two comparisons bound the answer. In the drag-only model a paddling stroke does not have to suppress the lift that its
flipper would produce, which favours it; in the lift model it has to, and the constrained search is difficult, which penalizes
it. Together they indicate that in this model a drag-based stroke needs between about two and ten times the power of a
lift-based stroke at the same speed, and that it cannot reach the speed at which lift-propelled gaits still cruise.

A complementary check replays the 27 feasible optima of the drag-only model with lift switched on. For 21 of them the lift
thrust share exceeds 0.1 (median 1.26), so the same leg motions are mainly lift-propelled once lift exists, and 26 no longer
stay straight and level. Whether a stroke is drag-based or lift-based is therefore a property of the physics as much as of the
motion: a flipper that crosses the water at an angle produces lift whether the model includes it or not. Paddling strokes found
in a drag-only model cannot be transferred unchanged to a real flipper.

\section{Summary}
\label{sec:mj-summary}

Within this model the answer to RQ1 is clear in direction. When lift is part of the physics and not suppressed, the optimizer selects lift-propelled strokes in all 101 feasible optima (22 in this chapter, the others in the coordination study of \cref{app:coord}). At equal speed they need one half to one third of the
power of the best strokes of a drag-only model. When drag-based strokes are forced within the lift model, the few that are
feasible (3 of 30 searches) need four to ten times the power of the lift-propelled strokes and do not reach
\SI{0.30}{\metre\per\second}; in this model a drag-based stroke therefore costs at least about twice and up to ten times as much. Whether a stroke is drag- or lift-based depends on the physics:
most paddling strokes of the drag-only model become lift-propelled when lift is added. Which coordination of the legs is best
is a separate question, which depends mainly on the attitude limits; it is reported in \cref{app:coord}. Three caveats limit these statements to the model: the coefficients and the lift
slope are not calibrated, the flat-plate lift has no aspect-ratio correction although the flipper has
$\mathit{AR}\approx1.4$, and the searches are not fully converged, particularly the constrained search for drag-based strokes. \Cref{ch:cfd,ch:results} therefore test the
central claim, that lift-based flapping beats drag-based paddling, with resolved flow simulations of the real leg.
